% ECONOMICS_PROBLEM: {"id": "L13-76", "title": "Algorithmic pricing conduct", "jel_code": "L13", "source_id": "OP-0398"}
% !TeX program = xelatex
\documentclass[11pt,a4paper]{article}
\usepackage[margin=23mm]{geometry}
\usepackage{fontspec}
\setmainfont{Latin Modern Roman}
\usepackage{amsmath,amssymb,mathtools}
\usepackage{xcolor,enumitem,booktabs,longtable,array,xurl,hyperref}
\definecolor{muted}{HTML}{53636C}
\hypersetup{colorlinks=true,urlcolor=blue,linkcolor=blue}
\setlength{\parindent}{0pt}
\setlength{\parskip}{5pt}
\newcommand{\R}{\mathbb R}
\newcommand{\E}{\mathbb E}
\newcommand{\N}{\mathbb N}
\newcommand{\ind}{\mathbf 1}
\newcommand{\Var}{\operatorname{Var}}
\newcommand{\Cov}{\operatorname{Cov}}
\newcommand{\supp}{\operatorname{supp}}
\DeclareMathOperator*{\argmin}{arg\,min}
\DeclareMathOperator*{\argmax}{arg\,max}
\newcommand{\entrymeta}[3]{{\sffamily\small\color{muted}\textbf{\texttt{#1}}\quad|\quad #2\par #3\par}}
\newcommand{\Problemref}[1]{\texttt{#1}}
\newcommand{\JELref}[2]{\texttt{#2} (JEL \texttt{#1})}
\begin{document}
\section*{Algorithmic pricing conduct}
\textbf{Problem ID:} \texttt{L13-76}\quad \textbf{JEL:} \texttt{L13}\par
% BEGIN ECONOMICS STATEMENT
\label{problem:OP-0398}
{\sffamily\small\color{muted}Original subject group: Industrial organization and competition\par}
\label{definitions:problem:30007680}
\textbf{Audit classification:} Explicit published open question. The source explicitly asks about pricing-algorithm welfare. The original phrase survives unilateral replacement did not define profits or a deviation criterion; a quantity-weighted mean also needed a nonzero denominator.
\subsection*{Definitions and precise research target}
Consider one market with \(N\geq2\) incumbent firms and an infinite sequence of periods; evaluate its first \(T\) periods for increasing positive integer horizons \(T\). Each active firm sets a price in a common compact interval, faces a specified demand function, and has observed production costs. Pricing programs learn from past public prices, quantities and their own profits; they exchange no direct messages. Potential entrants face a specified entry-cost distribution and decide when to enter. For an algorithm class \(\mathcal A\), define \(\bar p(a)\) as the quantity-weighted mean price generated by algorithm profile \(a\in\mathcal A^N\), averaged over learning randomness and demand shocks. Let \(p^B\) be the corresponding price benchmark from the same environment with independently optimal one-period pricing and identical entry opportunities and a specified equilibrium-selection rule; entrant algorithms are fixed separately in both comparisons. For a prespecified economically material threshold \(\varepsilon>0\), identify conditions under which \(\bar p(a)-p^B\geq\varepsilon\) persists as \(T\) increases and survives unilateral algorithm replacement. Include perturbations of demand, exploration rates and entrant behavior, and distinguish reciprocal punishment from unilateral pricing commitment. Report which conditions depend on algorithm architecture or market stationarity. This proposed dynamic specification extends a research agenda about sophisticated pricing; it is not a claim that a particular learning program necessarily colludes or that elevated prices prove communication.

\textbf{Definitions, assumptions and mathematical qualifications.} Fix one infinite-history environment and examine its first \(T\) periods, rather than changing the game arbitrarily with \(T\). Each program is a measurable map from its available history and random seed to a price; entrants' policies and entry law are input data. Define \(\bar p_T(a)=E[\sum_{tj}p_{jt}q_{jt}]/E[\sum_{tj}q_{jt}]\) with denominator bounded below by \(q_{\min}T>0\). Specify a static-pricing benchmark \(p_T^B\) using the same shocks and entry technology. Persistence means \(\liminf_{T\to\infty}(\bar p_T(a)-p_T^B)\geq\varepsilon\). To formalize survival of replacement, define \(U_{j,T}(a)=T^{-1}E[\sum_{t\leq T}(p_{jt}q_{jt}-C_j(q_{jt}))]\) and require \(\limsup_T\{U_{j,T}(a'_j,a_{-j})-U_{j,T}(a)\}\leq0\) for every allowed replacement \(a'_j\). An elevated price alone is not an equilibrium or a communication test.
\subsection*{Variants and assumption changes}
\begin{enumerate}
\item \textbf{Editorial, fixed learners.} For specified learning updates, exploration schedules and initialization distributions, estimate or prove price and profit limits; distinguish algorithm performance from stability to replacement.
\item \textbf{Editorial, program equilibrium.} Let firms choose from a finite or compact program menu and impose the profit-deviation condition. Compare stationary demand with a declared class of time-varying demand and entry shocks.
\end{enumerate}
\subsection*{References and source audit}
\textbf{Checked source.} 24 August 2024 author draft, Section 6, printed p. 37, fourth research direction. Full primary text retrieved. \url{https://web.stanford.edu/~ayurukog/trendsincompetition.pdf}.
\begin{enumerate}\raggedright
\item\hypertarget{definitions:source:30007680:1}{} Carl Shapiro; Ali Yurukoglu. 2024. \emph{Trends in Competition in the United States: What Does the Evidence Show?}. 24 August 2024 author draft, Section 6, printed p. 37, fourth research direction.
\url{https://web.stanford.edu/~ayurukog/trendsincompetition.pdf}.
DOI: \url{https://doi.org/10.3386/w32762}.
\end{enumerate}

\subsection*{Common conventions: Source mathematical conventions}

These conventions apply to each entry unless explicitly replaced.

\textbf{Models and identification.} A primitive model \(m\) specifies all probability laws, feasible choices, information, equilibrium and measurement rules used by its target. The admissible class \(\mathcal M\) is fixed independently of outcomes. If \(P_O(m)\) is its observable probability law and \(\tau(m)\) the target, its identified set at observable law \(P\) is
\[
 \mathcal I_\tau(P)=\{\tau(m):m\in\mathcal M,\ P_O(m)=P\}.
\]
Point identification means that this set is a singleton. An empty set signals incompatibility of data and assumptions. Coordinate bounds need not describe the joint set. Bounds are sharp only if every retained target is attainable or arbitrarily approachable by compatible models. Finite bounds require explicit restrictions on unobserved quantities; unrestricted confounding, hidden wealth, extrapolation or arbitrary equilibrium selection can yield uninformative sets.

\textbf{Causal interventions.} A potential outcome \(Y_i(p)\) is the outcome under a complete policy regime \(p\), including timing, financing, information and market spillovers. The target population and units are fixed before treatment. With interference, use the complete assignment vector or a justified exposure mapping; individual no-interference notation alone cannot identify an economy-wide intervention. A difference of mean potential outcomes does not require the cross-world joint distribution. Conditioning on post-treatment migration, survival, enrollment or employment changes the estimand and needs separate justification.

\textbf{Statistical statements.} An estimator, confidence set, forecast or test specifies a sampling law, model class and loss/coverage criterion. Uniform \((1-\alpha)\)-coverage means \(\inf_{m\in\mathcal M}\Pr_m\{\tau(m)\in C_n\}\geq1-\alpha\), or an explicitly asymptotic version. The whole parameter space trivially has coverage, so informativeness requires a stated width, risk or power objective. A forecasting improvement is relative to a named benchmark, information set and evaluation population.

\textbf{Equilibrium and optimization.} An equilibrium comprises feasible plans, optimality, beliefs where needed, budget constraints and all market/resource clearing. The primitives or a stated benchmark define each correspondence \(\mathcal E(p;m)\). Nonemptiness is assumed or proved, never inferred just from compact policy sets. A selected-equilibrium target uses a declared selection rule; a robust target quantifies over all permitted equilibria. Use suprema or infima unless attainment is established. An \(\varepsilon\)-optimal policy is feasible and within \(\varepsilon\) of the finite optimum value. Report an empty feasible set separately.

\textbf{Welfare and accounting.} Normative utility, interpersonal normalization, social weights, numeraire, discounting and the use of public receipts are primitives. Behavioral utility need not coincide with the welfare criterion. Household welfare includes ownership income and rents once; adding profits again to compensating variation generally double-counts them. A monetary equivalent is defined by an explicit transfer and price convention, with monotonicity and range conditions. Average effects, mechanism contributions, transfers and welfare losses are different targets. Infinite sums require convergence and dynamic feasible plans require terminal or no-Ponzi conditions.

\textbf{Source versions.} A working-paper date, revision date and journal publication year are distinguished. A DOI for a journal version is not automatically the DOI of an earlier manuscript. Locators refer to the stated version and printed pages where specified. A metadata-only check does not verify a claim about a full-text conclusion. Every entry's source subsection narrows the provenance claim accordingly.
% END ECONOMICS STATEMENT
\end{document}
